% !TEX program = xelatex
% 数学分析第三章：导数与微分

\documentclass[12pt, a4paper]{article}

% ---------- 基础宏包 ----------
\usepackage{fontspec}
\usepackage{xeCJK}
\setmainfont{Times New Roman}
\setCJKmainfont{SimSun}

\usepackage{amsmath, amssymb, amsthm}
\usepackage{geometry}
\geometry{left=3cm, right=3cm, top=2.5cm, bottom=2.5cm}

% ---------- 定理环境 ----------
\newtheorem{definition}{定义}[section]
\newtheorem{theorem}{定理}[section]
\newtheorem{corollary}{推论}[section]
\newtheorem{lemma}{引理}[section]
\newtheorem{example}{例}[section]

% ---------- 自定义命令 ----------
\newcommand{\R}{\mathbb{R}}
\newcommand{\dif}{\mathrm{d}} % 微分符号 upright d

% ---------- 文档开始 ----------
\begin{document}

\title{数学分析 III：导数与微分}
\author{inshergorohl}
\date{Flux 1.3 Reference}
\maketitle

\section{导数的概念}
\subsection{定义}

\begin{definition}[导数]
设函数 $y = f(x)$ 在 $x_0$ 的某邻域内有定义。若极限
\[
f'(x_0) = \lim_{\Delta x \to 0} \frac{f(x_0 + \Delta x) - f(x_0)}{\Delta x}
\]
存在，则称 $f$ 在 $x_0$ 处\textbf{可导}，并称该极限为 $f$ 在 $x_0$ 处的\textbf{导数}。
\end{definition}

\begin{definition}[左导数与右导数]
\[
f'_-(x_0) = \lim_{x \to x_0^-} \frac{f(x) - f(x_0)}{x - x_0}, \quad
f'_+(x_0) = \lim_{x \to x_0^+} \frac{f(x) - f(x_0)}{x - x_0}.
\]
\end{definition}

\begin{theorem}[可导与连续的关系]
若 $f(x)$ 在 $x_0$ 处可导，则 $f(x)$ 在 $x_0$ 处必连续。反之不成立（如 $y=|x|$ 在 $x=0$）。
\end{theorem}

\section{求导法则}

\subsection{基本公式}
\begin{enumerate}
    \item $(C)' = 0$
    \item $(x^\alpha)' = \alpha x^{\alpha-1}$
    \item $(\sin x)' = \cos x$, \quad $(\cos x)' = -\sin x$
    \item $(e^x)' = e^x$, \quad $(\ln x)' = \frac{1}{x}$
\end{enumerate}

\subsection{复合函数求导（链式法则）}
\begin{theorem}
设 $y = f(u), u = g(x)$ 均可导，则复合函数 $y = f(g(x))$ 可导，且
\[
\frac{\dif y}{\dif x} = \frac{\dif y}{\dif u} \cdot \frac{\dif u}{\dif x}.
\]
\end{theorem}

\section{微分中值定理}

\begin{theorem}[Fermat 引理]
若函数 $f$ 在 $x_0$ 处取极值且可导，则 $f'(x_0) = 0$。
\end{theorem}

\begin{theorem}[Rolle 中值定理]
设 $f \in C[a,b]$，且在 $(a,b)$ 内可导，且 $f(a)=f(b)$。
则 $\exists \xi \in (a,b)$，使得 $f'(\xi) = 0$。
\end{theorem}

\begin{theorem}[Lagrange 中值定理]
设 $f \in C[a,b]$，且在 $(a,b)$ 内可导。
则 $\exists \xi \in (a,b)$，使得
\[
f(b) - f(a) = f'(\xi)(b-a).
\]
\end{theorem}

\begin{corollary}[导数与单调性]
若 $f'(x) > 0$ 在 $(a,b)$ 内恒成立，则 $f(x)$ 在 $[a,b]$ 上严格单调增。
\end{corollary}

\section*{附录：导数速查表}
\[
\begin{array}{|c|c|}
\hline
f(x) & f'(x) \\ \hline
x^n & nx^{n-1} \\ \hline
\sin x & \cos x \\ \hline
\cos x & -\sin x \\ \hline
e^x & e^x \\ \hline
\ln x & 1/x \\ \hline
\end{array}
\]

\end{document}
